% Préambule commun aux cours de la formation C++.
% Inclus par chaque cours.tex via % Préambule commun aux cours de la formation C++.
% Inclus par chaque cours.tex via % Préambule commun aux cours de la formation C++.
% Inclus par chaque cours.tex via % Préambule commun aux cours de la formation C++.
% Inclus par chaque cours.tex via \input{preambule}.

\usepackage[utf8]{inputenc}
\usepackage[T1]{fontenc}
\usepackage[french]{babel}
\usepackage{lmodern}
\usepackage{microtype}

\usepackage{geometry}
\geometry{a4paper, margin=2.5cm}

\usepackage{xcolor}
\usepackage{amsmath, amssymb}
\usepackage{enumitem}
\usepackage{hyperref}
\hypersetup{
    colorlinks=true,
    linkcolor=blue!60!black,
    urlcolor=blue!60!black,
    citecolor=blue!60!black,
}

% --- Coloration syntaxique du C++ via listings -----------------------------
\usepackage{listings}

\definecolor{codebg}{rgb}{0.97,0.97,0.95}
\definecolor{codekw}{rgb}{0.10,0.10,0.60}
\definecolor{codecomment}{rgb}{0.20,0.50,0.20}
\definecolor{codestring}{rgb}{0.65,0.15,0.15}
\definecolor{coderule}{rgb}{0.80,0.80,0.80}

% Permet les caractères accentués (UTF-8) à l'intérieur des blocs de code.
\lstset{
    literate=%
        {é}{{\'e}}1 {è}{{\`e}}1 {ê}{{\^e}}1 {ë}{{\"e}}1
        {à}{{\`a}}1 {â}{{\^a}}1 {ä}{{\"a}}1
        {î}{{\^i}}1 {ï}{{\"i}}1
        {ô}{{\^o}}1 {ö}{{\"o}}1
        {ù}{{\`u}}1 {û}{{\^u}}1 {ü}{{\"u}}1
        {ç}{{\c c}}1 {É}{{\'E}}1 {È}{{\`E}}1 {À}{{\`A}}1
}

\lstdefinestyle{cpp}{
    language=C++,
    backgroundcolor=\color{codebg},
    basicstyle=\ttfamily\small,
    keywordstyle=\color{codekw}\bfseries,
    commentstyle=\color{codecomment}\itshape,
    stringstyle=\color{codestring},
    showstringspaces=false,
    numbers=left,
    numberstyle=\tiny\color{gray},
    numbersep=8pt,
    frame=single,
    rulecolor=\color{coderule},
    breaklines=true,
    tabsize=4,
    morekeywords={constexpr,noexcept,override,final,nullptr,auto,
        std,string,vector,size_t,uint64_t,int64_t,optional,namespace}
}
\lstset{style=cpp}

% Raccourci pour du code en ligne, en police machine à écrire.
\newcommand{\code}[1]{\texttt{#1}}

% Boîte « À retenir »
\usepackage{tcolorbox}
\newtcolorbox{aretenir}[1][]{
    colback=yellow!8, colframe=orange!70!black,
    title=#1, fonttitle=\bfseries, sharp corners, boxrule=0.5pt
}
\newtcolorbox{piege}[1][Piège]{
    colback=red!5, colframe=red!60!black,
    title=#1, fonttitle=\bfseries, sharp corners, boxrule=0.5pt
}
.

\usepackage[utf8]{inputenc}
\usepackage[T1]{fontenc}
\usepackage[french]{babel}
\usepackage{lmodern}
\usepackage{microtype}

\usepackage{geometry}
\geometry{a4paper, margin=2.5cm}

\usepackage{xcolor}
\usepackage{amsmath, amssymb}
\usepackage{enumitem}
\usepackage{hyperref}
\hypersetup{
    colorlinks=true,
    linkcolor=blue!60!black,
    urlcolor=blue!60!black,
    citecolor=blue!60!black,
}

% --- Coloration syntaxique du C++ via listings -----------------------------
\usepackage{listings}

\definecolor{codebg}{rgb}{0.97,0.97,0.95}
\definecolor{codekw}{rgb}{0.10,0.10,0.60}
\definecolor{codecomment}{rgb}{0.20,0.50,0.20}
\definecolor{codestring}{rgb}{0.65,0.15,0.15}
\definecolor{coderule}{rgb}{0.80,0.80,0.80}

% Permet les caractères accentués (UTF-8) à l'intérieur des blocs de code.
\lstset{
    literate=%
        {é}{{\'e}}1 {è}{{\`e}}1 {ê}{{\^e}}1 {ë}{{\"e}}1
        {à}{{\`a}}1 {â}{{\^a}}1 {ä}{{\"a}}1
        {î}{{\^i}}1 {ï}{{\"i}}1
        {ô}{{\^o}}1 {ö}{{\"o}}1
        {ù}{{\`u}}1 {û}{{\^u}}1 {ü}{{\"u}}1
        {ç}{{\c c}}1 {É}{{\'E}}1 {È}{{\`E}}1 {À}{{\`A}}1
}

\lstdefinestyle{cpp}{
    language=C++,
    backgroundcolor=\color{codebg},
    basicstyle=\ttfamily\small,
    keywordstyle=\color{codekw}\bfseries,
    commentstyle=\color{codecomment}\itshape,
    stringstyle=\color{codestring},
    showstringspaces=false,
    numbers=left,
    numberstyle=\tiny\color{gray},
    numbersep=8pt,
    frame=single,
    rulecolor=\color{coderule},
    breaklines=true,
    tabsize=4,
    morekeywords={constexpr,noexcept,override,final,nullptr,auto,
        std,string,vector,size_t,uint64_t,int64_t,optional,namespace}
}
\lstset{style=cpp}

% Raccourci pour du code en ligne, en police machine à écrire.
\newcommand{\code}[1]{\texttt{#1}}

% Boîte « À retenir »
\usepackage{tcolorbox}
\newtcolorbox{aretenir}[1][]{
    colback=yellow!8, colframe=orange!70!black,
    title=#1, fonttitle=\bfseries, sharp corners, boxrule=0.5pt
}
\newtcolorbox{piege}[1][Piège]{
    colback=red!5, colframe=red!60!black,
    title=#1, fonttitle=\bfseries, sharp corners, boxrule=0.5pt
}
.

\usepackage[utf8]{inputenc}
\usepackage[T1]{fontenc}
\usepackage[french]{babel}
\usepackage{lmodern}
\usepackage{microtype}

\usepackage{geometry}
\geometry{a4paper, margin=2.5cm}

\usepackage{xcolor}
\usepackage{amsmath, amssymb}
\usepackage{enumitem}
\usepackage{hyperref}
\hypersetup{
    colorlinks=true,
    linkcolor=blue!60!black,
    urlcolor=blue!60!black,
    citecolor=blue!60!black,
}

% --- Coloration syntaxique du C++ via listings -----------------------------
\usepackage{listings}

\definecolor{codebg}{rgb}{0.97,0.97,0.95}
\definecolor{codekw}{rgb}{0.10,0.10,0.60}
\definecolor{codecomment}{rgb}{0.20,0.50,0.20}
\definecolor{codestring}{rgb}{0.65,0.15,0.15}
\definecolor{coderule}{rgb}{0.80,0.80,0.80}

% Permet les caractères accentués (UTF-8) à l'intérieur des blocs de code.
\lstset{
    literate=%
        {é}{{\'e}}1 {è}{{\`e}}1 {ê}{{\^e}}1 {ë}{{\"e}}1
        {à}{{\`a}}1 {â}{{\^a}}1 {ä}{{\"a}}1
        {î}{{\^i}}1 {ï}{{\"i}}1
        {ô}{{\^o}}1 {ö}{{\"o}}1
        {ù}{{\`u}}1 {û}{{\^u}}1 {ü}{{\"u}}1
        {ç}{{\c c}}1 {É}{{\'E}}1 {È}{{\`E}}1 {À}{{\`A}}1
}

\lstdefinestyle{cpp}{
    language=C++,
    backgroundcolor=\color{codebg},
    basicstyle=\ttfamily\small,
    keywordstyle=\color{codekw}\bfseries,
    commentstyle=\color{codecomment}\itshape,
    stringstyle=\color{codestring},
    showstringspaces=false,
    numbers=left,
    numberstyle=\tiny\color{gray},
    numbersep=8pt,
    frame=single,
    rulecolor=\color{coderule},
    breaklines=true,
    tabsize=4,
    morekeywords={constexpr,noexcept,override,final,nullptr,auto,
        std,string,vector,size_t,uint64_t,int64_t,optional,namespace}
}
\lstset{style=cpp}

% Raccourci pour du code en ligne, en police machine à écrire.
\newcommand{\code}[1]{\texttt{#1}}

% Boîte « À retenir »
\usepackage{tcolorbox}
\newtcolorbox{aretenir}[1][]{
    colback=yellow!8, colframe=orange!70!black,
    title=#1, fonttitle=\bfseries, sharp corners, boxrule=0.5pt
}
\newtcolorbox{piege}[1][Piège]{
    colback=red!5, colframe=red!60!black,
    title=#1, fonttitle=\bfseries, sharp corners, boxrule=0.5pt
}
.

\usepackage[utf8]{inputenc}
\usepackage[T1]{fontenc}
\usepackage[french]{babel}
\usepackage{lmodern}
\usepackage{microtype}

\usepackage{geometry}
\geometry{a4paper, margin=2.5cm}

\usepackage{xcolor}
\usepackage{amsmath, amssymb}
\usepackage{enumitem}
\usepackage{hyperref}
\hypersetup{
    colorlinks=true,
    linkcolor=blue!60!black,
    urlcolor=blue!60!black,
    citecolor=blue!60!black,
}

% --- Coloration syntaxique du C++ via listings -----------------------------
\usepackage{listings}

\definecolor{codebg}{rgb}{0.97,0.97,0.95}
\definecolor{codekw}{rgb}{0.10,0.10,0.60}
\definecolor{codecomment}{rgb}{0.20,0.50,0.20}
\definecolor{codestring}{rgb}{0.65,0.15,0.15}
\definecolor{coderule}{rgb}{0.80,0.80,0.80}

% Permet les caractères accentués (UTF-8) à l'intérieur des blocs de code.
\lstset{
    literate=%
        {é}{{\'e}}1 {è}{{\`e}}1 {ê}{{\^e}}1 {ë}{{\"e}}1
        {à}{{\`a}}1 {â}{{\^a}}1 {ä}{{\"a}}1
        {î}{{\^i}}1 {ï}{{\"i}}1
        {ô}{{\^o}}1 {ö}{{\"o}}1
        {ù}{{\`u}}1 {û}{{\^u}}1 {ü}{{\"u}}1
        {ç}{{\c c}}1 {É}{{\'E}}1 {È}{{\`E}}1 {À}{{\`A}}1
}

\lstdefinestyle{cpp}{
    language=C++,
    backgroundcolor=\color{codebg},
    basicstyle=\ttfamily\small,
    keywordstyle=\color{codekw}\bfseries,
    commentstyle=\color{codecomment}\itshape,
    stringstyle=\color{codestring},
    showstringspaces=false,
    numbers=left,
    numberstyle=\tiny\color{gray},
    numbersep=8pt,
    frame=single,
    rulecolor=\color{coderule},
    breaklines=true,
    tabsize=4,
    morekeywords={constexpr,noexcept,override,final,nullptr,auto,
        std,string,vector,size_t,uint64_t,int64_t,optional,namespace}
}
\lstset{style=cpp}

% Raccourci pour du code en ligne, en police machine à écrire.
\newcommand{\code}[1]{\texttt{#1}}

% Boîte « À retenir »
\usepackage{tcolorbox}
\newtcolorbox{aretenir}[1][]{
    colback=yellow!8, colframe=orange!70!black,
    title=#1, fonttitle=\bfseries, sharp corners, boxrule=0.5pt
}
\newtcolorbox{piege}[1][Piège]{
    colback=red!5, colframe=red!60!black,
    title=#1, fonttitle=\bfseries, sharp corners, boxrule=0.5pt
}
